% TonalWM — OJSP FULL VERSION: the growing complete study.
% The 4-page ICASSP submission (paper/icassp2027.tex) is CUT DOWN from this file.
% Plan per SPEC section 9, unchanged: 4p ICASSP first; the full study goes to OJSP.
% Every measurement lands here; only a subset survives the 4-page cut.
% TonalWM — full study
% ==========================================================================
% INTEGRITY (SPEC §7): every number is traceable to a ledgered artifact under
% results/. NO number is estimated or placeholdered. "TBD (run required)" marks
% a measurement that has not been made. Numbers carry a [sN] tag naming the run.
%
% FORMAT: written for spconf.sty (official ICASSP Paper Kit). spconf is not
% redistributable, so this file falls back to IEEEtran when it is absent — the
% fallback exists only so the draft compiles; drop spconf.sty in the directory
% and it is used automatically. No content depends on which is active.
% ==========================================================================
\documentclass[conference]{IEEEtran}
\newif\ifspconf
\IfFileExists{spconf.sty}{\spconftrue}{\spconffalse}
\ifspconf \usepackage{spconf} \fi
% IEEEtran fallback so the draft always compiles. spconf is not redistributable;
% drop spconf.sty beside this file and it is used instead, with no content change.
\ifspconf\else
  \newcommand{\name}[1]{\author{\IEEEauthorblockN{#1}}}
  \newcommand{\address}[1]{}
  \newcommand{\ninept}{}
  \newenvironment{keywords}{\begin{IEEEkeywords}}{\end{IEEEkeywords}}
\fi
\usepackage[T1]{fontenc}
\usepackage{lmodern}
\usepackage{amsmath,amssymb,graphicx,booktabs,array,tabularx}
\usepackage[hidelinks]{hyperref}
\usepackage{url}
\graphicspath{{../results/figures/}}   % figures are ledgered results/ artifacts

\newcommand{\tkr}{\mathrm{TKR}}
\newcommand{\ikr}{\mathrm{IKR}}
\newcommand{\ecyc}{\varepsilon_{\mathrm{cyc}}}
\newcommand{\ppl}{\delta_{\mathrm{PPL}}}
\newcommand{\TBD}[1]{\textbf{[TBD (run required): #1]}}
\newcommand{\vs}{vs.\@\ }

\title{Does a Music Transformer Know What Key It Is In?\\
Causal Interventions on an Emergent Tonal World Model}

\name{Anonymous ICASSP 2027 submission}
\address{Paper ID: TBD}

\begin{document}
\ninept
\maketitle

\begin{abstract}
Music generation models produce convincingly tonal music. Whether they do so by
chaining surface statistics or by maintaining an internal representation of
musical key that generation \emph{causally uses} has not been tested. We test it
directly: a decoder-only Transformer trained on symbolic music from a
functional-harmony grammar whose vocabulary contains absolute pitch and time
only---no key, chord, or degree tokens---under a pre-registered protocol with
frozen decision rules. Our central result is causal: replacing the key component
of the residual stream at a single middle layer makes the model continue in the
injected key at $5.0\times$ the rate of rank- and norm-matched random edits
(guarded strict $\tkr$ $0.378$ \vs $0.075$; significant for $12/12$ target keys
after Holm correction, replicated across seeds), while staying inside a
musicality budget frozen before any edit ran---reaching $58\%$ of the behavioral
ceiling obtained by transposing the prompt itself. A capacity sweep sharpens what
this means. Next-token accuracy is flat from $0.5$M to $86$M parameters, yet at
$0.5$M the key is not decodable beyond the input surface at all, and the same
edit moves the key \emph{more} often than at $26$M while destroying the music in
$90\%$ of cases. Capacity does not buy the key; it buys a key that can be moved
without breaking everything else. Finally we leave our own models entirely. On real Bach
chorales a \emph{public} $128$M checkpoint trained on real music---architecturally
identical to one of ours---carries a key state that beats the pitch surface where ours
does not, and editing that state redirects the music it composes at $5.7\times$ the rate
of matched random edits ($11/12$ targets, Holm-corrected) with \emph{no} measurable cost
to musicality. It also resists: the keys the edit cannot reach are exactly the keys real
music rarely uses (Spearman $+0.91$). Editing a music Transformer's internal key state
changes what key it composes in---in a model people actually use, not only in ours.
\end{abstract}

\begin{keywords}
music generation, world models, interpretability, causal intervention, tonality
\end{keywords}

\section{Introduction}
\label{sec:intro}
Systems that compose music now write convincingly tonal pieces: their outputs
settle into a home key, build tension away from it, and return. For listeners
this is much of what makes generated music sound like music; for creators it
determines whether such systems can be steered.

How they achieve this is largely unexamined. Existing analyses report that
musical attributes can be \emph{read out} of a model's activations by trained
classifiers~\cite{syntheory2024}. Reading, however, is not using: a signal can be
decodable from a system's state without the system consulting it when it
acts---a gap documented repeatedly in interpretability
research~\cite{hewitt2019control}. No study has tested whether a music model
\emph{uses} an internal key representation: whether intervening on it, while
leaving the input untouched, changes the key of the music the model goes on to
generate.

Sequence models can build causal models of the latent state behind their data. A
model trained only to predict Othello moves tracks the board those moves imply,
and editing that internal board changes which moves it treats as
legal~\cite{li2023othello,nanda2023emergent}. Analogous causal evidence exists
for chess, mazes, and RL environments~\cite{karvonen2024chess,maze2024,iris2026},
but not for music.

We close that gap. We choose key---tonic and mode---as the state variable for
three reasons. First, unlike any candidate in the game domains, key is defined
\emph{externally} to the model, by centuries of theory and by psychological
measurement~\cite{krumhansl1982}, so the probed state is not an artifact of our
own labeling. Second, among musical state candidates it is the most demanding:
metric position is recoverable from time tokens, and harmonic function
presupposes a key, whereas key must be \emph{computed} by integrating pitch
evidence over time. Third, key admits a natural analogue of Othello's legal-move
rate---diatonic scale membership---so causal success is measurable without a
learned judge.

Our contributions:
\begin{enumerate}\itemsep1pt
\item A controlled testbed for world-model analysis in music: a leak-free
tokenizer, a functional-harmony corpus with exact per-token key labels, and a
pre-registered protocol, released as a benchmark (Sec.~\ref{sec:method}).
\item Existence and geometry: key is decodable beyond strong input-side
baselines in all six trained models, the twelve tonic directions self-organize
into a circle of fifths, and---a pre-registered negative result---transposition
augmentation does not make the representation a cyclic group
(Sec.~\ref{sec:resultsA}).
\item Causality: single-layer subspace edits redirect the composed key for
$12/12$ targets within a frozen musicality budget, and a capacity sweep shows
that what capacity supplies is not the key but the ability to move it
surgically---dissociating three things the term ``world model'' usually conflates
(Secs.~\ref{sec:resultsB}--\ref{sec:capacity}).
\item A matched-architecture control on real music: with corpus, protocol, and decision
rule held fixed, a public real-trained model clears the input-surface baseline where our
synthetic-trained model of the same shape does not, isolating the training data as the
variable that decides whether a key state is worth more than the notes themselves
(Sec.~\ref{sec:wild}).
\item The same causal test on that public model: editing its key state redirects the key
it composes in ($5.7\times$ the matched control, $11/12$ targets, $100\%$ within the
frozen musicality budget), so the world-model claim is not confined to models we built.
Its failures are informative---it resists exactly the keys real music rarely uses
(Sec.~\ref{sec:wildcausal}).
\end{enumerate}

\section{Relation to Prior Work}
\label{sec:prior}
Prior work divides along two axes: domain and evidence level. Across domains,
world-model analyses have progressed from readability to causal intervention in
board games~\cite{li2023othello,nanda2023emergent,mcgrath2022chess,karvonen2024chess},
mazes~\cite{maze2024}, and RL environments~\cite{iris2026}. In music, evidence has
stayed at the readability level: probing studies decode theory concepts from
generation models~\cite{syntheory2024}, and steering studies control attributes
without testing a state variable~\cite{smitin2024}. We occupy the empty
cell---causal, state-level evidence in music---for a state whose transitions
(modulations) themselves follow a grammar, unlike the game domains.

\section{Method}
\label{sec:method}

\subsection{Problem statement and terminology}
The input is a symbolic music sequence; the latent state is the running key
$\kappa(t)$; the question is whether the generator's continuation causally
depends on its internal estimate of $\kappa$. Table~\ref{tab:terms} fixes the
terms. The pre-registration (supplementary) specifies every decision rule,
threshold, and control, all frozen before any experiment ran; deviations are
logged and negative results reported.

\begin{table}[t]
\setlength{\tabcolsep}{4pt}\renewcommand{\arraystretch}{1.3}
\centering\small
\caption{Terminology. Decision rules (DR-$*$) were frozen pre-experiment.}
\label{tab:terms}
\begin{tabularx}{\linewidth}{@{}l>{\raggedright\arraybackslash}X@{}}
\toprule
$\kappa(t)$ & Key state: tonic $(\mathbb{Z}_{12})\times{}$mode at token $t$\\
$\kappa^{*}$ & Donor key injected by an edit\\
D-SYN & Synthetic corpus with per-token ground-truth $\kappa(t)$\\
M-CTRL & Studied model\\
M-REF & Frozen reference model; used only for the quality guard\\
Edit & $h \leftarrow h - P_{V}h + P_{V}\mu_{\kappa^{*}}$ in the residual stream\\
$\tkr$ & Rate at which the continuation's estimated key $=\kappa^{*}$\\
$\ikr_{x}$ & Fraction of continuation pitches diatonic to key $x$\\
$\ppl$ & M-REF perplexity budget, frozen pre-intervention\\
\bottomrule
\end{tabularx}
\end{table}

\begin{figure}[t]
\centering
\includegraphics[width=\columnwidth]{fig_framework.pdf}
\caption{The claim in one example (real model output, fixed sampling seed [s0]).
The same $8$-bar B$\flat$-major prompt, continued (a)~clean and (b)~with the key
component of layer-$4$ activations replaced by the E-major class mean from bar~$9$
on. \textbf{The input does not change}; the continuation's estimated key follows
the injected internal state. Shading marks the scale of the key in force.}
\label{fig:overview}
\end{figure}

\subsection{Pre-registered hypotheses}
\label{sec:hyp}
The protocol, its thresholds and its decision rules were frozen before any experiment
ran. We restate the hypotheses in the pre-registration's own terms rather than
paraphrase them, since the decision rules below are only meaningful against the claims
they were written for.

\begin{itemize}\itemsep2pt
\item[\textbf{H1}] (\emph{existence}) Key is decodable from internal activations, and
more accurately than from input-surface baselines (recent-window pitch-class histogram
/ Krumhansl--Schmuckler); the gap widens in ambiguous regions. \emph{Falsified if} no
layer's selectivity-corrected probe beats the input baseline.
\item[\textbf{H2}] (\emph{encoding form}) The key representation is
transposition-equivariant: for each transposition $T_k$ there is a linear map $R_k$ in
representation space, and $\{R_k\}$ has cyclic-group structure ($R_k \approx R_1^{\,k}$).
\item[\textbf{H2b}] A model trained with transposition augmentation (R-Aug) shows
\emph{stronger} equivariance than one without (R-NoAug)---a manipulation of the training
condition. H2b is a sub-claim of H2, which is why no ``H2a'' appears: the
pre-registration numbered the parent claim H2 and its manipulation H2b.
\item[\textbf{H3}] (\emph{causality; the main claim}) A targeted edit of the key
subspace moves the key of the continuation to the edit target. The effect significantly
exceeds a random subspace edit \textbf{matched in rank and norm}, and stays inside the
musicality guard's budget. \emph{Falsified if} $\tkr$ is indistinguishable from the
control edit, or is achievable only together with quality collapse.
\item[\textbf{H4}] (\emph{persistence and reassertion}) After a \emph{one-shot} edit the
state either (a)~persists, (b)~is reasserted toward the original key by context, or
(c)~transitions via a pivot-like passage. The pre-registration assigned H4 to
Phase~C; its sub-question H4a (does an installed value persist at all, and through
what carrier?) became load-bearing for the paper's framing and was brought forward.
Sec.~\ref{sec:h4a} reports it; the remainder of H4 (pivot structure, annotation)
stays in Phase~C.
\item[\textbf{H5}] (\emph{localization}) The layers at which an edit is causally
effective are confined to a limited band.
\end{itemize}

\noindent H3 asks for a control matched in \emph{rank and norm}, and ``norm'' admits two
readings; we checked both rather than assert the convenient one. The basis reading is
satisfied by construction: $V$ and $V_{\mathrm{K1}}$ are both orthonormal $512\times24$,
so $\lVert V\rVert_F = \lVert V_{\mathrm{K1}}\rVert_F = \sqrt{24}$ identically. The
stricter reading---that the two edits should \emph{perturb the residual stream equally}
---is not imposed by construction, since we never rescale the injected component, so we
measured it on real activations during real generation at the tested condition
($12$ targets, all edited positions). The edit applies
$\lVert\Delta\rVert = 21.9$ against K1's $22.1$: a ratio of $0.99$, or $0.160$ \vs
$0.162$ of $\lVert h\rVert = 136.5$, with the \emph{control} marginally the larger. The
match is observed rather than enforced, but it is observed, and it is conservative.

That the two agree so closely is itself informative. Both edits project
$\mu_{\kappa^{*}}\!-\!h$ onto a rank-$24$ subspace of $\mathbb{R}^{512}$; a random one
captures $\sqrt{24/512}\approx0.217$ of that vector's energy, and the key subspace is
measured to capture $0.215$---no more. The difference between $\mu_{\kappa^{*}}$ and a
given activation is dominated by variation that has nothing to do with key, and the key
subspace holds no privileged share of it. So the $5\times$ gap reported in
Sec.~\ref{sec:results} cannot be a magnitude effect: the two interventions are the same
size, and only their direction differs.

\subsection{Data and models}
To make ground truth exact rather than estimated, all training data are
synthetic. \textbf{D-SYN.} A functional-harmony state machine
(T\,$\to$\,S\,$\to$\,D\,$\to$\,T with substitutions, SATB voicing, optional
diatonic melody) generates $200$k$/10$k$/10$k train/val/test pieces of
$278$--$508$ tokens with $0$--$3$ modulations each (pivot-chord, direct, or
sequential; fifths-distance stratified), giving every token a ground-truth key
label. The three types are planned in equal thirds, but two of them degrade in the
token stream: a pivot with no shared diatonic triad emits a direct jump, and a
sequential shift over an odd fifths-distance cannot be halved into two equal steps.
The realized corpus is therefore $18.6\%$ pivot, $16.5\%$ two-step sequential, and
$64.9\%$ single jumps---a fact the surface baselines benefit from, not our probe. We
also note that our ``sequential'' is a documented simplification (two equal-interval
shifts one bar apart), not a full melodic sequence. Every token still carries an exact key
label. \textbf{Tokenizer.} To prevent state leakage the vocabulary's $124$ types are
\texttt{BAR}, \texttt{POS}, absolute-MIDI \texttt{PITCH} and \texttt{DUR} ($121$ types)
together with the three structural symbols \texttt{PAD}/\texttt{BOS}/\texttt{EOS}, which
are key-invariant. Nothing names a key, chord or degree, and a unit test enforces this.
\textbf{Models.} M-CTRL is a GPT-2 style decoder ($8$ layers, $8$ heads,
$d\,{=}\,512$, context $512$, ${\sim}25$M parameters), trained $12$k steps under
two regimes---with (R-Aug) and without (R-NoAug) uniform $12$-fold transposition
augmentation---$\times\,3$ seeds. M-REF shares the architecture but uses a
disjoint seed \emph{and} data split, and serves only as the musicality guard,
avoiding circular evaluation. All models pass a pre-frozen quality gate
(validation next-token top-$1$ $0.878$--$0.881$ \vs a $0.434$ constant-predictor
rate; unconditional in-key ratio $0.964$--$0.977$ \vs a $0.608$ threshold).

\subsection{Phase A: probing and geometry}
To separate internal computation from input transcription, every probe result is
reported against three control families. Linear (multinomial logistic) probes map
residual-stream activations $h_{\ell}(t)$ to $\kappa(t)$ ($24$ classes) with
sequence-level splits. \textbf{C1} permutes key identities per sequence---a
control task in the sense of~\cite{hewitt2019control}; \textbf{C2} is an untrained
model; \textbf{C3} are input-side baselines at the same positions (pitch-class
histogram logistic regression, and Krumhansl--Schmuckler (KS) template matching;
windows $W\in\{8,16,32,64\}$). DR-H1: the probe supports H1 iff its
selectivity-corrected $F_1$ exceeds the best C3 baseline with a sequence-level BCa
$95\%$ CI excluding zero in some layer. Ambiguity is stratified by KS posterior
entropy ($W{=}16$). For geometry, orthogonal Procrustes maps $R_k$ are fit between
activations of pieces and their $k$-semitone transpositions, and a cyclicity error
tests group structure:
\begin{equation}
\ecyc \;=\; \frac{1}{10}\sum_{k=2}^{11}
\frac{\lVert R_k - R_1^{\,k}\rVert_F}{\lVert R_k\rVert_F},
\label{eq:cyc}
\end{equation}
normalising each term by $\lVert R_k\rVert_F$ before averaging, as pre-registered. The
sum starts at $k{=}2$ because $R_1 - R_1^{1}$ is identically zero and would only dilute
the mean, and stops at $k{=}11$ because $T_{12}$ is the identity; $R_1$ defines the
powers. DR-H2b: $\ecyc(\text{R-Aug}) < \ecyc(\text{R-NoAug})$ (one-sided Wilcoxon
across seeds).

\subsection{Phase B: causal intervention}
\label{sec:methodB}
To test \emph{use} rather than presence, we edit the state and score the behavior.
From bar~$9$ of a key-stable $8$-bar prompt, a sustained edit replaces the key
component at one layer, $h \leftarrow h - P_V h + P_V\mu_{\kappa^{*}}$, with $V$ an
orthonormal key subspace and $\mu_{\kappa^{*}}$ the class-conditional mean of the
donor key. ``Sustained'' is meant literally and asymmetrically, and the asymmetry is
worth stating because it is easy to assume otherwise: the edit is applied at
\emph{every} token position from $t^{*}$ (the bar-$9$ boundary) onward---\texttt{BAR},
\texttt{POS}, \texttt{PITCH} and \texttt{DUR} alike---and re-applied at every generation
step, whereas the \emph{probe} is trained and scored at sampled positions, uniformly
over token types (the predict-pitch reading convention enters only with the
public-model study, Sec.~\ref{sec:wild}). (What a \emph{one-shot} installation leaves behind---and through what
carrier---is Sec.~\ref{sec:h4a}.) We write to the whole stream and read from
one position in it. Restricting the edit to pitch-choosing positions is a natural
variant and we have not run it. Subspaces: \textbf{V-PROBE} (probe row space, rank $24$),
\textbf{V-MEAN} (class-mean span), \textbf{V-DAS} (interchange-trained, ranks
$8/24$). Controls: \textbf{K1} rank- and norm-matched random subspace
(Sec.~\ref{sec:hyp}); \textbf{K2} sham
edit, required bit-identical to the clean run (passed, $100/100$ prompts);
\textbf{K3} a different layer's subspace; \textbf{K4} behavioral reference,
transposing the prompt itself. Continuations ($16$ bars; $12$ major targets
$\times\,100$ prompts $\times\,8$ layers; sampling noise paired with a clean twin)
are scored by KS key estimation, $\ikr$, and a grammar guard. $\tkr$ is reported
\emph{strict} (the estimated key equals $\kappa^{*}$ exactly) and \emph{tolerant}. The
pre-registration defines tolerant as the disjunction of three relations, not fifths
adjacency alone: fifths distance ${\leq}\,1$ \emph{with the same mode}, the relative
major/minor (mode differs, tonics $3$ or $9$ semitones apart), or the parallel key
(mode differs, same tonic). Every headline number in this paper is strict; tolerant
exists to bracket the estimator's known confusions (Sec.~\ref{sec:disc}), not to soften
the claim. The budget $\ppl\,{=}\,0.613$ ($90$th percentile of
the M-REF perplexity rise across $7{,}989$ natural modulations in validation data)
was fixed \emph{before} any edit ran; edits exceeding it do not count as successes.
DR-H3 requires $\tkr(\text{edit})>\tkr(\text{K1})$ with Holm-corrected Wilcoxon
$p<.05$ in ${\geq}\,8/12$ targets at the best in-budget condition.

\section{Results}
\label{sec:results}
Intervention numbers are from the R-Aug seed-$0$ sweep [s0] unless noted; probing
holds for all six M-CTRL models, and the seed-$1$ sweep replicates every Phase-B
verdict (Sec.~\ref{sec:repl}).

\subsection{H1: the key is decodable beyond the input surface}
\label{sec:resultsA}
The key is \emph{computed} by the model, not transcribed from the input: the
corrected probe--baseline margin excludes zero in layers $1$--$7$ of every trained
model (peak $+0.105$). It fails only at L$0$, and fails there in the informative
direction: the margin is significantly \emph{negative} ($-0.20$, CI
$[-0.211,-0.182]$), i.e.\ at the raw embedding the input surface outperforms the
model's own representation of it. Table~\ref{tab:probe} gives
the numbers at the layer the edit acts on.

\subsubsection{The history-length confound, and what survived it}
\label{sec:extD}
The comparison above has an asymmetry that could have explained it entirely: the
probe reads $h_\ell(t)$, which attends over the full prefix ($\le512$ tokens,
${\sim}170$ notes), while C3 sees a $W$-token window with $W\le64$ (${\sim}21$
notes)---and the best C3 sat exactly at $W{=}64$, the \emph{edge} of the
pre-registered sweep, with its $F_1$ still rising steeply ($0.686$ at $W{=}32$).
Both the headline margin and the ambiguity contrast could then be read as
``the probe sees more history,'' not ``the model computes a state.''

We therefore reran C3 at the same positions with the sweep extended to the full
prefix ($W\in\{96,128,192,256,512\}$; $512$ exceeds the longest piece) and with two
post-hoc variants that only strengthen the baseline: exponentially recency-decayed
full-history histograms (half-life $\lambda\in\{32,\ldots,256\}$ tokens) and a
concatenation of short and full windows $[W16\,|\,W512]$. History alone does not
close the gap. The plain-window curve turns over \emph{inside} the extension, at
$W{=}96$ ($0.801$), and falls monotonically to $0.752$ at the full prefix: with
$0$--$3$ modulations per piece, evidence from before a modulation poisons a
histogram that cannot down-weight it, which is precisely the estimation problem an
internal state solves. The strongest surface overall is the concatenation, at
$0.824$; the decayed variant peaks at $0.814$ ($\lambda{=}32$). Against that
strongest baseline the corrected margin at L$4$ is $+0.071$ (BCa $95\%$ CI
$[0.056,0.081]$), and every DR-H1 verdict in this paper---$6/6$ main-line support,
both $0.5$M failures, both $3.4$M and both $86$M passes---is unchanged, though the
$3.4$M margins thin to $+0.015$ $[0.004,0.023]$ and $+0.012$ $[0.001,0.021]$.

One registered sub-claim does \emph{not} survive. Against the pre-registered
baseline the probe's advantage concentrated where KS estimation is uncertain
($0.936$ \vs $0.795$ in the top ambiguity quartile); against the strengthened
baseline the advantage is \emph{uniform}---$+0.097$ in the top quartile, $+0.097$
elsewhere. The apparent ambiguity-specific edge was an artifact of the $W{=}64$
cap, which is weakest exactly where a short window is uncertain. We report H1's
``the gap widens where the key is ambiguous'' as not supported in its strengthened
form, and Fig.~\ref{fig:ambiguity} should be read with that caveat.

\begin{figure}[t]
\centering
\includegraphics[width=0.8\columnwidth]{fig_ambiguity.pdf}
\caption{A cautionary result. Against the \emph{pre-registered} baseline (shown),
the probe's advantage concentrates in the top ambiguity quartile ($0.94$ \vs
$0.79$)---which reads as ``the state earns its keep where the surface fails.''
Against the strengthened baseline of Sec.~\ref{sec:extD} the advantage is uniform
across strata ($+0.097$ in both): the concentration was an artifact of the
$W{=}64$ window, itself weakest where a short window is uncertain.}
\label{fig:ambiguity}
\end{figure}

\begin{table}[t]
\setlength{\tabcolsep}{4pt}\renewcommand{\arraystretch}{1.15}
\centering\small
\caption{Key decoding at L$4$ (R-Aug seed~$0$), $24$-class macro-$F_1$ on held-out
sequences; DR-H1 supported in $6/6$ models. Note the two C1 rows measure different
things: a control probe trained on permuted labels barely learns \emph{them}
($0.040$, near the $1/24$ chance rate---the control task is genuinely unlearnable),
and predicts the \emph{true} key at $0.032$---the selectivity floor the margin
subtracts. The margin reconciles exactly:
$0.927-0.032-0.790=0.105$.}
\label{tab:probe}
\begin{tabularx}{\linewidth}{@{}Xr@{}}
\toprule
\textbf{Method / baseline} & \textbf{macro-$F_1$}\\
\midrule
Linear probe on $h_4(t)$ & $0.927$\\
\quad MLP probe (secondary) & $0.932$\\
Best input baseline (C3: PC-hist LR, $W{=}64$) & $0.790$\\
Untrained model (C2) & $0.243$\\
C1 control probe, on its own permuted labels & $0.040$\\
C1 control probe, on the true key (selectivity floor) & $0.032$\\
\midrule
Corrected margin, BCa $95\%$ CI $[0.091,\,0.114]$ & $+0.105$\\
\bottomrule
\end{tabularx}
\end{table}

\subsection{H2: fifths geometry; equivariant but not a group}
The representation carries music-theoretic structure, but not the strictest form we
pre-registered. Projected to two principal components, the twelve major-tonic
directions order as C--G--D--A--E--B--F$\sharp$--$\cdots$ (Fig.~\ref{fig:fifths}):
circle-of-fifths neighbour agreement $1.00$ \vs $0.00$ for the chromatic ordering,
at every probed layer (exploratory). Transposition acts as a generalizing linear
map: Procrustes $R_k$ fit on half the sequences aligns held-out activations with
$0.05$--$0.11$ relative error. The family is not cyclic, however: $\ecyc\approx0.98$
\vs ${\approx}\,1.41$ for random orthogonal maps [Eq.~\eqref{eq:cyc}]. Our
pre-registered comparison is negative: augmentation does not improve cyclicity
($3$ seeds per regime; R-Aug $0.983$ \vs R-NoAug $0.981$; one-sided $p=0.83$), so
DR-H2b is \emph{not} supported. The representation is transposition-equivariant in
effect but not organized as a cyclic group, and the key diversity already present
in the data appears to suffice.

\begin{figure}[t]
\centering
\includegraphics[width=0.62\columnwidth]{fig_fifths_geometry.pdf}
\caption{The twelve major-tonic probe directions (L$4$, first two PCs)
self-organize into a circle of fifths. No key, chord, or degree token exists in
the training vocabulary.}
\label{fig:fifths}
\end{figure}

\subsection{H3: editing the key state changes the composed key}
\label{sec:resultsB}
Editing the internal key redirects generation. At the best in-budget condition
(V-PROBE, L$4$), $\tkr(\text{edit})-\tkr(\text{K1})>0$ holds for $12/12$ targets
(Holm-corrected $p<10^{-4}$ each; rank-biserial $r=0.64$--$1.00$) against a
pre-registered bar of ${\geq}\,8/12$: DR-H3 is supported. Fig.~\ref{fig:bars} gives
all conditions. Pitch content triangulates the same conclusion without the
estimator: $\ikr_{\kappa^{*}}=0.937$ \vs $\ikr_{\kappa_{\mathrm{src}}}=0.646$ after
editing. Success falls gently with target distance on the circle of fifths ($0.395$
at distance~$1$ to $0.280$ at the tritone), while random edits collapse to zero once
the identity target is excluded. The edit reaches $58\%$ of the behavioral ceiling
(K4), quantifying an \emph{actionability gap}: the state is readable in a rank-$24$
subspace of one layer, yet a single-layer edit realizes only part of the behavioral
effect. The fraction depends on a convention we should state. K4's perplexity excess is
measured against the \emph{untransposed} clean twin, so it scores the transposition
rather than any damage, and guarding K4 is not well defined. We therefore divide the
edit's guarded rate ($0.378$) by K4's raw rate ($0.652$), giving $58\%$. Applying the
guard to both sides gives $0.378/0.614 = 62\%$ and to neither $0.429/0.652 = 66\%$; we
report the smallest of the three, which is the one that flatters us least. Two subsidiary findings: the \emph{optimized} subspace (V-DAS) underperforms
the probe's readout directions (Sec.~\ref{sec:disc}), and V-MEAN edits destroy output
quality at deep layers (guard pass ${\leq}\,5\%$ at L$4$--L$5$)---the guard is
load-bearing, a point Sec.~\ref{sec:capacity} makes decisive.

The blanket write itself---every token position, one type-averaged
$\mu_{\kappa^{*}}$---is a design one should suspect of damaging the music
needlessly, so we tested it: the same condition with the write restricted by token
type. It is vindicated, and the refutation is informative. Restricting to PITCH
positions is causally \emph{null} (guarded $0.077$, at the K1 floor) while leaving
the music essentially untouched (guard pass $0.974$); restricting to POS${+}$PITCH
recovers $0.156$; the \emph{complement}---BAR and DUR positions only---recovers
$0.261$, more than POS${+}$PITCH from fewer positions; and the two halves sum
approximately to the whole ($0.156{+}0.261$ \vs $0.429$ raw). Guard rates do not
improve under any mask, and the perplexity cost per unit of guarded $\tkr$ is
\emph{lowest} for the blanket write: the damage tracks the effect, not the write's
breadth. The causally potent component of the key state thus sits mostly at
bar-line and duration positions, not at the pitch emissions themselves---a third
dissociation alongside Sec.~\ref{sec:h5}'s, whose mechanism we leave open: probe
$F_1$ by token type is flat, per-type edit displacement is comparable, and
attention at the reading layers concentrates on PITCH sources, so none of the
three obvious accounts (readability, edit magnitude, attention anchoring)
explains it.

\begin{figure}[t]
\centering
\includegraphics[width=0.95\columnwidth]{fig_intervention_bars.pdf}
\caption{Guarded strict $\tkr$ across conditions [s0]: the continuation lands
exactly on the injected key \emph{and} stays within the frozen musicality budget.
Chance $=1/12$. Bars: prompt-level bootstrap $95\%$ CI.}
\label{fig:bars}
\end{figure}

\subsection{H5: readability is diffuse; causal efficacy is sharp}
\label{sec:h5}
The layer profile of the edit-minus-control effect is single-peaked
($0.03,0.04,0.19,0.26,\mathbf{0.30},0.26,0.23,0.23$ for L$0$--L$7$); the
best-\vs-median bootstrap CI $[0.038,0.109]$ excludes zero, so DR-H5 is supported.
The two profiles in Fig.~\ref{fig:layers} are \emph{not} the same shape, and the
difference is the finding. The probe saturates across L$2$--L$6$ (macro-$F_1$
$0.920$--$0.929$, all within $0.01$ of its maximum), so the key is legible from most
of the network; yet a single-layer edit moves the output only near L$4$, and barely
at all at L$0$--L$1$. Where a state is legible is a weak guide to where it is
used---a caution for interpretability work that localizes a computation by probing
accuracy alone.

\begin{figure}[t]
\centering
\includegraphics[width=\columnwidth]{fig_layer_profile.pdf}
\caption{(a)~Probe macro-$F_1$ per layer, all six M-CTRL models (thin lines; R-Aug
seed~$0$ highlighted): the key is legible across a broad plateau. (b)~Causal effect
[s0]: guarded strict $\tkr$ of the V-PROBE edit \vs the K1 matched random control
(bands: prompt-level bootstrap $95\%$ CI). Only a narrow band around L$4$ moves the
output.}
\label{fig:layers}
\end{figure}

\subsection{H4a: the variable is re-estimated, not stored}
\label{sec:h4a}
Everything above edits under a \emph{sustained} clamp, and a clamp cannot witness
the ``install a value'' semantics our framing borrows: perhaps the model merely
tolerates being steered, and maintains nothing. So we installed the value
\emph{once} and let go. With no key--value cache the past is recomputed at every
step, so ``once'' means: the edit is pinned to the positions of the first generated
bar (each row's window closes at its own next bar line) and no later position is
ever touched; later positions may attend to the pinned bar-$9$ state and to the
notes it caused. Sustained, clean, and a K1 one-shot bracket it; all conditions
share sampling randomness, and the one-shot's first bar is \emph{identical} to the
sustained run's by construction ($\ikr_{\kappa^{*}}=0.805$ in both, an internal
validity check).

Released, the effect does not decay---and it does not complete
(Fig.~\ref{fig:persistence}a). $\ikr_{\kappa^{*}}$ drops once ($0.805\to0.685$)
and then holds a flat plateau ($0.65$--$0.70$) for thirteen bars, $+0.12$ above
the K1 floor of $0.55$ (average scale overlap between random major keys). The
music neither returns to the source key nor settles in the target: per-bar KS
names the target in only ${\sim}9\%$ of bars (the clamp climbs to $28\%$), and
under a strict last-four-bars metric the continuation returns to the source in
$31\%$ of runs---\emph{less} often than after a matched random perturbation
($45\%$), which itself matches the clean run's own $47\%$: a random one-bar
perturbation derails nothing; the installed bar durably shifts the tonal center
of mass without finishing the modulation.

Then the control that decides what ``persists'' means
(Fig.~\ref{fig:persistence}b): splice the one-shot run's bar-$9$ \emph{tokens}
after the clean prompt and continue with no activation edit anywhere. The splice
reproduces the plateau in full ($0.687$ \vs $0.671$ across bars $2$--$14$; state
pinning adds $-0.016$, i.e.\ nothing). \textbf{The carrier is the music, not the
pinned state}: a value installed in the activations persists exactly insofar as
it writes itself into the notes, which then stand as ordinary input evidence.
The key variable this paper's probes read and this paper's edits move is a
continuously re-estimated statistic of the evidence---causally consulted at
every step (H3), never latched (H4a). Our framing must therefore not rest the
``set a variable'' distinction on persistence-after-release, and does not: it
rests on the remove-and-install semantics, the matched controls, and the
dissociations of Secs.~\ref{sec:resultsB}--\ref{sec:h5}.

\begin{figure}[t]
\centering
\includegraphics[width=\columnwidth]{fig_persistence.pdf}
\caption{H4a. (a)~Released after a one-bar installation, the pitch-content shift
neither decays nor completes: a flat plateau $+0.12$ over the K1 floor for $13$
bars. (b)~Splicing the installed bar's \emph{tokens} with no activation edit
reproduces the plateau; pinning the state adds nothing. The variable is
re-estimated from evidence, not stored.}
\label{fig:persistence}
\end{figure}

\subsection{Replication, and what capacity actually buys}
\label{sec:capacity}
\label{sec:repl}
\textbf{Second seed.} An independently trained R-Aug seed-$1$ model reproduces every
Phase-B verdict: DR-H3 supported with V-PROBE significant in $12/12$ targets
($\tkr$ $0.23$--$0.41$ \vs K1 $0.02$--$0.10$), and DR-H5 supported (CI
$[0.122,0.181]$). Its profile peaks at L$2$ rather than L$4$: the \emph{localization}
replicates, its exact depth does not, and we report both profiles rather than a
pooled one.

\textbf{Capacity.} Training the same data and protocol at $0.5$M (L$2$, $d{=}128$),
$3.4$M (L$4$, $d{=}256$), $26$M (main line), and $86$M (L$12$, $d{=}768$) parameters
(counted from the checkpoints) dissociates three things the term ``world model''
usually conflates (Fig.~\ref{fig:emergence}). \emph{Behavior} is flat: next-token
top-$1$ is $0.868$--$0.879$ across a $170\times$ capacity range. \emph{Readability} is
not: at $0.5$M the probe peaks at $F_1$ $0.63$--$0.66$, \emph{below} the input baseline
($0.79$), and DR-H1 is not supported in either seed---a behaviorally equivalent model
that has no key state at all, which is the surface-statistics account made flesh, and
which our pre-registered test correctly refuses to credit. (All size verdicts are
unchanged against the strengthened baseline of Sec.~\ref{sec:extD}; the $3.4$M
margins thin to $+0.015$ $[0.004,0.023]$ and $+0.012$ $[0.001,0.021]$---the state
first appears just barely above the strongest surface.) \emph{Manipulability} is a
third thing again. At $0.5$M the edit changes the key \emph{more} often than at $26$M
in raw terms ($0.55$ \vs $0.43$) and passes the musicality guard only $10\%$ of the
time (\vs $80\%$ at $26$M): with two layers, ``editing the key subspace'' is a blunt
overwrite of the output, not the movement of a separable state. Reporting guarded
$\tkr$ alone would hide this; reporting raw $\tkr$ alone would credit a surface model
with a world model. \textbf{Capacity does not buy the key---it buys a key that can be
moved without breaking everything else}, and the pre-registered guard is precisely the
instrument that tells the two apart.

Manipulability does not simply keep improving, however. It peaks near $26$M and then
\emph{falls}: the $86$M model reaches guarded $\tkr$ $0.18$, about half the $26$M value.
We suspected an artifact of our own making---each size model was edited at its
probe-$F_1$ peak, a heuristic this paper itself refutes---so we swept all twelve layers
of the $86$M model. The hypothesis was wrong: its causal peak (L$6$, $0.206$) is
statistically indistinguishable from the layer we had already used (L$3$, $0.202$).
The decline is real. A natural reading, which we did not test and do not claim, is that
greater depth leaves more downstream layers able to re-derive the key from the surviving
context, so a single-layer edit is more easily repaired---the actionability gap widening
with scale rather than closing. (Layer sweep: one seed.)

\begin{figure}[t]
\centering
\includegraphics[width=0.95\columnwidth]{fig_emergence.pdf}
\caption{Behavioral equivalence, representational dissociation. Across a $170\times$
capacity range next-token accuracy is flat, but the key readout is not: below
${\sim}6$M the probe cannot beat the input baseline and the pre-registered DR-H1
margin is negative. Dots: seeds.}
\label{fig:emergence}
\end{figure}

\subsection{Real music, and a public model: what generalizes}
\label{sec:wild}
Two questions remain for a synthetic-data study: does the key readout survive real
music, and does any of this describe a model people actually use? We answer both with
one matched comparison, on $300$ Bach chorales whose \emph{local} key at every note is
given by published Roman-numeral analyses.\footnote{Scores: \texttt{bach-370-chorales}
(CC BY-NC-SA); analyses: \emph{When in Rome} (CC BY-SA). Neither is redistributed.}

\textbf{Our model does not clear the surface on real music.} Probed on the chorales, our
$86$M model reaches macro-$F_1$ $0.534$ against a selectivity floor of $0.035$---the key
subspace is real and transfers, it is not an artifact of our grammar---but the best input
baseline (KS, $W{=}64$) reaches $0.451$, and the corrected margin, $+0.034$ (CI
$[-0.021,\,0.103]$), does not exclude zero. DR-H1 is \emph{not} supported out of
distribution.

\textbf{A public model trained on real music does clear it.} The Anticipatory Music
Transformer~\cite{thickstun2023anticipatory} (Apache-2.0; Lakh MIDI, MetaMIDI, FMA, and
$450$k commercial records) is, at its \texttt{small} size, \emph{architecturally
identical} to our $86$M model: $12$ layers, $d{=}768$, and $85.1$M non-embedding
parameters against our $85.2$M. Only the totals differ ($128$M \vs $86$M), and only
because its vocabulary holds $55$k tokens where ours holds $124$---so the stacks being
compared are the same size, and what differs is what they were trained on. Its vocabulary
is (arrival-time, duration, note) triples---leak-free in exactly the sense ours is. We
probed it on the same chorales, at the same positions, with the same controls and the
same decision rule; the two C3 baselines land at $0.451$ and $0.454$, confirming the
setups are matched. It reaches macro-$F_1$ $0.690$ (L$10$) with a corrected margin of
$+0.196$ (CI $[0.132,\,0.260]$): \textbf{DR-H1 supported}. With architecture, corpus,
protocol, and decision rule held fixed, \emph{the training data is the single free
variable}, and it is what decides whether a key state beats the pitch surface. The layer
profiles say the same thing from the other side: our synthetic model's readout peaks
early (L$2$) and decays to $0.17$ by L$11$ on this out-of-distribution input, while the
real-trained model's dips mid-network and \emph{builds} to its peak at L$10$.

One caution, and it is a large one. The probe must be read where the model is about to
\emph{choose} a pitch (the token before a note). Read at the note itself---after the
pitch is committed, when the model is predicting the next onset---the same public model,
on the same chorales, with the same probe, controls and rule, produces a different
picture and a different verdict. Its layer profile inverts from a U-shape peaking deep
($F_1$ $0.618$ at L$0$, $0.357$ at L$4$, $0.690$ at L$10$) to a monotone decay ($0.545$,
$0.227$, $0.093$), and the corrected margin collapses from $+0.196$ (CI $[0.132,0.260]$)
to $+0.042$ (CI $[-0.034,0.092]$), which does not exclude zero: read at the note,
\textbf{DR-H1 is not supported} on this model. Nothing changed but the position. Had we
probed where the pitch has already been emitted, we would have concluded that a public
music Transformer's key state does not beat the pitch surface and that it fades with
depth---and we would have published it. Where one probes is not a free parameter, and we
report the negative convention alongside the positive one rather than only the one that
worked.

\subsection{The public model \emph{uses} its key state}
\label{sec:wildcausal}
Readability is not use, and the whole point of this paper is that the difference is
testable. We therefore put the public model through the same intervention its synthetic
counterpart faced. The subspace is the row space of the probe we fit on it; the guard is
a \emph{different} public checkpoint (music-medium refereeing music-small: different
weights, so not circular), with $\ppl$ frozen at $0.849$ nats---the $90$th percentile of
that reference's perplexity rise across $1{,}180$ \emph{natural} modulations in the
chorales. Because this paper shows the probe peak is not the causal peak, the edit layer
cannot be inherited from the probe: it is searched on $20$ chorale prefixes and then
judged on $60$ prefixes it never saw. The sham edit is bit-identical to the clean run
($20/20$), so any movement is the edit's.

We state the freeze order precisely, because it differs from the main experiment's. For
our own models the budget was frozen before \emph{any} edit ran. Here it was frozen after
the $20$-prefix layer scan and before the held-out evaluation---the order the ledger
records. The scan never consults the budget (it ranks layers on unguarded
$\tkr(\text{edit})-\tkr(\text{K1})$), and every number we report comes from the
evaluation that followed the freeze; but the scan's raw outcomes were visible when the
budget's parameters were fixed, so this is a weaker guarantee than the pre-registered one
and we do not claim otherwise.

\textbf{DR-H3 is supported} (Table~\ref{tab:wildcausal}). At L$8$, on held-out prompts,
the edit lands the continuation on the injected key at $5.7\times$ the rate of the
matched random control ($0.365$ \vs $0.064$), significant in $11/12$ targets after Holm
correction. $\ikr$ crosses: the continuation is \emph{more} diatonic to the injected key
than to the prompt's own ($0.824$ \vs $0.785$). And the edit costs the music nothing
measurable---$100\%$ of edited continuations sit inside the frozen budget (median excess
$+0.125$ nats against a budget of $0.849$), where our own model managed $87\%$. A key
state learned from real music is \emph{more} cleanly separable, not less. One
consequence deserves stating plainly, since the guard does discriminating work elsewhere
in this paper: here it does none. Both the edit and K1 pass it at $100\%$, so guarded and
raw $\tkr$ coincide exactly ($0.365$ and $0.064$) and the word ``guarded'' in
Table~\ref{tab:wildcausal} filters nothing. That is the finding, not a gap in it---but
the reader should not picture the budget separating survivors from casualties on this
model, because there are no casualties.

Two findings here are only visible on a real model. First, the causal peak (L$8$) is
again \emph{not} the probe peak (L$10$): the diffuse-readout/sharp-effect dissociation of
Sec.~\ref{sec:h5} is not an artifact of synthetic data. Second, \textbf{the model has a
key prior and resists being pushed out of it}. The single target that fails outright is
F$\sharp$ major ($\tkr$ $0.000$), and across all twelve targets success is almost
perfectly ordered by how common that tonic is in the chorale corpus. We count key
frequency two ways over the same $300$ chorales, both from the human Roman-numeral
analyses. Weighting each key by how much music is actually spent in it (note-level local
labels, $n{=}138{,}249$) gives Spearman $+0.91$ against $\tkr$; counting each chorale's
opening key once ($n{=}300$) gives $+0.75$. With only twelve tonics the asymptotic test
is unreliable, so we report one-sided Monte-Carlo permutation $p$ (direction predicted a
priori): $p<10^{-4}$ and $p=0.003$ respectively. On the note-weighted count, D at
$13.8\%$ of the corpus reaches $\tkr$ $0.600$ and G at $21.4\%$ reaches $0.550$, while
D$\flat$ at $0.2\%$ reaches $0.233$ and G$\flat$ at $2.2\%$ reaches zero.

Two limits of this analysis deserve naming. The chorale corpus is a \emph{proxy} for
which keys Western tonal music favours; it is not this model's training distribution,
which is Lakh MIDI, MetaMIDI and FMA~\cite{thickstun2023anticipatory}. The correlation
therefore establishes that the model has a key prior aligned with tonal practice, not
that the prior was inherited from any particular corpus---identifying its source would
require the training distribution itself. And G$\flat$ is the one target where the prior
argument and the estimator's known weakness coincide: it is both the rarest sharp-side
key and the enharmonic seam of the circle, so its zero should be read as the ordering's
endpoint, not as independent evidence. Our synthetic model, whose key distribution is
uniform by construction, could not have shown any of this.

\begin{table}[t]
\setlength{\tabcolsep}{4pt}\renewcommand{\arraystretch}{1.15}
\centering\small
\caption{Editing the key state of a \emph{public} model trained on real music
(music-small, $128$M, Apache-2.0), on $60$ held-out Bach chorale prefixes. The edit layer
was chosen on a disjoint set of prompts. Guard: a different public checkpoint, budget
frozen before any edit.}
\label{tab:wildcausal}
\begin{tabularx}{\linewidth}{@{}Xr@{}}
\toprule
Guarded strict $\tkr$, V-PROBE edit (L$8$) & $\mathbf{0.365}$\\
Guarded strict $\tkr$, K1 matched random control & $0.064$\\
\quad ratio & $5.7\times$\\
\midrule
Targets significant after Holm correction & $\mathbf{11/12}$\\
Guard pass rate (edit) & $100\%$\\
$\ikr$ to the injected key \vs to the prompt's key & $0.824$ \vs $0.785$\\
\bottomrule
\end{tabularx}
\end{table}

\section{Discussion and Limitations}
\label{sec:disc}
The three levels support the world-model reading \emph{within this training
condition}: a Transformer trained on notes alone maintains a mid-layer key variable
with music-theoretic geometry, and generation consults it. The gap to the behavioral
ceiling and the failure of edits at L$0$--L$1$ indicate a distributed, redundantly
re-estimated state rather than a single canonical register---a contrast with
Othello-GPT's near-complete editability. Sec.~\ref{sec:h4a} turns the natural
attribution into a measurement: music continually re-grounds tonality in the
immediately preceding notes, and accordingly an installed key value survives only
through the notes it writes---the state is re-estimated, never latched.

\begin{table}[t]
\centering\caption{K3 (layer-shuffled subspace) resolved by layer. K3 injects layer
$\ell{+}4$'s subspace at layer $\ell$; guarded strict $\tkr$, R-Aug seed~$0$. The
layer-pooled mean ($0.160$) averages two opposite regimes and describes no condition
that was run.}
\label{tab:k3}
\begin{tabular}{lccccc}
\toprule
inject at & L$0$ & L$1$ & L$2$ & L$3$ & L$4$ \\
(borrowed from) & (L$4$) & (L$5$) & (L$6$) & (L$7$) & (L$0$) \\
\midrule
K3            & $0.003$ & $0.274$ & $0.387$ & $0.313$ & $0.064$ \\
own edit      & $0.103$ & $0.118$ & $0.267$ & $0.337$ & $\mathbf{0.378}$ \\
K1            & $0.073$ & $0.074$ & $0.074$ & $0.078$ & $0.075$ \\
\bottomrule
\end{tabular}
\end{table}

K3 deserves its own paragraph, because its layer-pooled value misleads and we had
initially read it wrongly. Pooled over layers K3 is $0.160$, which sits invitingly
``between K1 and the true edit'' and reads as a graded partial effect. Resolved by layer
(Table~\ref{tab:k3}) it is nothing of the kind. At the tested layer L$4$, K3 borrows
L$0$'s subspace---the one layer whose probe fails to beat the input surface---and
accordingly does nothing ($0.064$, at or below K1's $0.075$). That is the clean negative
control the design intended, and it says the L$4$ effect requires L$4$'s own subspace.
Everywhere else the result inverts: subspaces borrowed from L$5$--L$7$ and injected at
L$1$--L$3$ reach $0.274$--$0.387$, and at L$2$ the borrowed subspace ($0.387$)
outperforms that layer's own edit ($0.267$) and even the L$4$ headline. The key subspace
is thus largely \emph{shared} across the middle stack rather than re-learned at each
depth; what depth governs is not which direction encodes the key but whether writing to
it changes the output. K3 is a valid negative control at L$4$ only, and we no longer read
its pooled mean as evidence of anything.

Five observations bound our claims. \emph{Mode:} the pre-registration (SPEC §4) fixed
the injected mode to major and designated minor targets a secondary condition. They have
not been run. Every causal claim in this paper is therefore a claim about major targets,
and whether a minor key state is equally manipulable---or whether the major/relative-minor
distinction is even carried in the same subspace---is open. \emph{V-DAS below V-PROBE:}
the optimized subspace may simply be undertrained relative to its search space; we report
it as observed and release its configuration, but do not claim optimization is inherently
inferior. \emph{Estimator:} the KS judge confuses fifths-neighbours, which
strict/tolerant reporting brackets and $\ikr$ triangulates. \emph{Capacity:} the sweep
varies depth and width together, so it locates a threshold, not its cause. \emph{External
validity:} Sec.~\ref{sec:wildcausal} carries both the readout and the causal result to a
public model trained on real music, so the world-model claim is no longer confined to our
own synthetic setting. What remains ours alone is the \emph{exactness} of the ground
truth: only the synthetic corpus gives a per-token key label that no estimator had to
guess. Whether edited continuations modulate
\emph{musically} (pivot-like transitions \vs jumps) is measured in the journal version.
All training data are synthetic; no artist recordings or copyrighted corpora are used
for training, and the evaluation corpora are used under their stated licenses.

\section{Conclusion}
\label{sec:concl}
We gave pre-registered causal evidence that a music Transformer tracks musical key as
an internal state variable that generation consults, and quantified how far a
single-layer edit goes toward full behavioral control. The capacity sweep separates
three properties usually bundled under ``world model''---the state being present,
being legible, and being separately manipulable---and shows they emerge at different
points. The testbed (exact state labels, leak-free vocabulary, frozen decision rules)
is released as a benchmark. Because music offers state variables that theory defines
and ears can verify, we expect it to be a useful substrate for asking what sequence
models know; our next question is whether models also internalize how keys are
\emph{supposed} to change.

\vspace{2pt}\noindent\footnotesize\textbf{Reproducibility.} Anonymized code, data
generator, pre-registration, and an audio demo page (clean \vs edited continuations):
\url{https://anonymous.4open.science/r/TBD}. \TBD{create before submission.}\normalsize

\begin{thebibliography}{9}\footnotesize
\bibitem{li2023othello} K.~Li, A.~K. Hopkins, D.~Bau, F.~Vi\'egas, H.~Pfister, and
M.~Wattenberg, ``Emergent world representations: Exploring a sequence model trained
on a synthetic task,'' in \emph{ICLR}, 2023.
\bibitem{nanda2023emergent} N.~Nanda, A.~Lee, and M.~Wattenberg, ``Emergent linear
representations in world models of self-supervised sequence models,'' in
\emph{BlackboxNLP}, 2023.
\bibitem{mcgrath2022chess} T.~McGrath et al., ``Acquisition of chess knowledge in
AlphaZero,'' \emph{PNAS}, vol.~119, no.~47, 2022.
\bibitem{karvonen2024chess} A.~Karvonen, ``Emergent world models and latent variable
estimation in chess-playing language models,'' in \emph{CoLM}, 2024.
\bibitem{maze2024} A.~F. Spies, W.~Edwards, M.~I. Ivanitskiy, et al., ``Transformers use
causal world models in maze-solving tasks,'' \emph{arXiv:2412.11867}, 2024.
\bibitem{iris2026} X.~Zhang, ``What do world models learn in RL? Probing latent
representations in learned environment simulators,'' \emph{arXiv:2603.21546}, 2026.
\bibitem{syntheory2024} M.~Wei, M.~Freeman, C.~Donahue, and C.~Sun, ``Do music generation
models encode music theory?,'' in \emph{ISMIR}, 2024.
\bibitem{smitin2024} I.~Prokopiou, P.~Vikatos, M.~Kaliakatsos-Papakostas, T.~Giannakopoulos,
and T.~Stafylakis, ``Latent space disentanglement via activation steering for
interpretable attribute control in symbolic music generation,'' in \emph{EUSIPCO}, 2026.
\bibitem{thickstun2023anticipatory} J.~Thickstun, D.~Hall, C.~Donahue, and P.~Liang,
``Anticipatory music transformer,'' \emph{TMLR}, 2024.
\bibitem{krumhansl1982} C.~L. Krumhansl and E.~J. Kessler, ``Tracing the dynamic
changes in the perceived tonal organization in a spatial representation of musical
keys,'' \emph{Psychological Review}, vol.~89, no.~4, pp.~334--368, 1982.
\bibitem{hewitt2019control} J.~Hewitt and P.~Liang, ``Designing and interpreting
probes with control tasks,'' in \emph{EMNLP}, 2019.
\end{thebibliography}

\end{document}
